\section{PVI (Problemas de Valor Inicial)}
Problemas de ecuaciones diferenciales en los que la variable independiente es el tiempo.
Las distintas aproximaciones no son soluciones más cercanas a la buscada, que reducen el error con iteraciones, sino que son soluciones correspondientes a avances temporales de paso h.
\subsection{Método Euler}
\begin{itemize}
    \item Error local de Euler es $O(h^2)$
    \item Error global $\approx O(h)$
    \item El método Euler tiene orden de precisión 1.
\end{itemize}
Forma recursiva: $u_{n+1}=u_n+h\ f(u_n,t_n)$
\subsection{Método de punto medio}
Avanzar con la pendiente evaluada en el punto medio del intervalo.\\
\underline{Forma recursiva}\\
\begin{enumerate}
\item $u_{n+\frac{1}{2}}=u_n+h\ f(u_n,t_n)$ Primer iteración por Euler para obetner $u_n+\frac{1}{2}$.
\item $u_{n+1}=u_n+h\ f(u_{n+\frac{1}{2}},t_{n+\frac{1}{2}})$
\end{enumerate}
\subsection{Método de Runge-Kutta}
El método Runge-Kutta es de orden 2.
\begin{enumerate}
    \item $q_1=h\ f(u_n,t_n)$ Toma la pendiente al inicio del intervalo.
    \item $q_1=h\ f(u_{n+q_1},t_{n+q_1})$ Toma la pendiente al final del intervalo.
    \item $u_{n+1}=u_n+(\frac{q_1+q_2}{2})$ Avanzar con el promedio de las pendientes
\end{enumerate}
\subsection{Método de Euler Implícito}
El método Runge-Kutta es de orden 1.
$u_{n+1}=u_n+h\ f(u_{n+1}, t_{n+1})$

